% ---------------------------------------------------------------
\section{Proposed Methodology}
\label{sec:method}

\subsection{Problem Formulation}
\label{sec:formulation}
A binary classifier $f:\mathbb{R}^{M}\rightarrow[0,1]$ on $M$ input features is
trained once and then frozen. It returns $p(x)=\Pr_f(Y{=}1\mid x)$, the prediction
$\hat{y}=\mathbb{I}\{p(x)\ge 0.5\}$ and the raw confidence
$s(x)=\max\{p(x),1-p(x)\}$. Before deployment, three disjoint labelled historical
partitions are drawn: a training set $\mathcal{D}_{\mathrm{tr}}$, a
correctness-calibration set $\mathcal{D}_{\mathrm{cal}}$ and a reference set
$\mathcal{D}_{\mathrm{ref}}$. The historical reference accuracy
$\Aref$ is the accuracy of $f$ on $\mathcal{D}_{\mathrm{ref}}$.
After deployment, unlabelled inputs arrive in batches
$B_t=\{x_{t,1},\dots,x_{t,n}\}$, $t=1,\dots,T$, of size $n$. The realised batch
accuracy $A_t=\frac{1}{n}\sum_{i}\mathbb{I}\{\hat{y}_{t,i}=y_{t,i}\}$ cannot be
computed when $B_t$ arrives, because its labels are released only after a delay of
$d\in\mathbb{N}$ batches. A batch is \emph{harmful} when
\begin{equation}
H_t=\mathbb{I}\{A_t<\Aref-\delta\},\qquad \delta=0.05.
\label{eq:harm}
\end{equation}
At every clock $t$ the monitor must report, from $B_t$ and the labels released so
far, (i) whether the input distribution has changed, (ii) which features carry the
change, (iii) an estimate $\Ahat_t$ of $A_t$ with an uncertainty statement, and
(iv) whether $H_t=1$ is likely. Because $P(X)$ can stay fixed while $P(Y\mid X)$
changes, (i) and (iv) are separate targets. The framework is required to keep them
separate and to show when (iv) cannot be supported by unlabelled evidence.

\subsection{Architecture Overview}
Fig.~\ref{fig:arch} shows the \idm{} pipeline. All artefacts are fitted on
historical data and then frozen. Two evidence channels run in parallel on each
batch: a label-free \emph{input channel} (detection and localization) and an
\emph{impact channel} (calibrated accuracy estimation and a harm-concern score).
A label-release broker is the only component that holds ground truth. It
reveals labels to a verification ledger strictly according to the delay schedule.
An aligned evidence view shows both channels and the verified outcomes on one
shared batch axis. It leaves a cell empty while its label is unreleased and never
imputes the value. The numbered badges 1--6 in Fig.~\ref{fig:arch} mark the stages
described below, and Table~\ref{tab:notation} in the appendix lists all symbols.

\begin{figure*}[t]
\centering
\resizebox{\textwidth}{!}{% Architecture diagram of IDM-Viz (TikZ; contains no data values).
\begin{tikzpicture}[
  font=\footnotesize,
  >={Stealth[length=2.1mm]},
  box/.style={draw=#1!80!black, line width=0.6pt, rounded corners=2.5pt, align=center, fill=#1!7, minimum height=10mm, inner sep=3pt, drop shadow={opacity=0.18, shadow xshift=0.5mm, shadow yshift=-0.5mm}},
  box/.default=c1,
  lane/.style={draw=#1!70, dashed, line width=0.6pt, rounded corners=5pt, inner sep=5pt, fill=#1!3},
  arr/.style={->, line width=0.7pt, draw=cink!75},
  lab/.style={font=\scriptsize\bfseries, text=#1!80!black},
  badge/.style={circle, fill=#1, text=white, font=\scriptsize\bfseries, inner sep=0.8pt, minimum size=4mm},
]
% --- frozen artefacts (top row)
\node[box=cgray, text width=27mm] (fit) {\textbf{Fit \& freeze}\\$f$, scaler $\phi$, PCA $W$\\on $\mathcal{D}_{\mathrm{tr}}$};
\node[box=cgray, text width=27mm, right=6mm of fit] (cal) {\textbf{Correctness map}\\isotonic $g$: $s\mapsto q$\\on $\mathcal{D}_{\mathrm{cal}}$};
\node[box=cgray, text width=27mm, right=6mm of cal] (ref) {\textbf{Reference}\\$\Aref$ on $\mathcal{D}_{\mathrm{ref}}$;\\KS sample $R$};
\node[box=cgray, text width=27mm, right=6mm of ref] (null) {\textbf{Null calibration}\\threshold $\tau_\Delta$ from\\stationary batches};

% --- stream
\node[box=c4, text width=22mm, below=13mm of fit.south west, anchor=north west, minimum height=36mm] (batch) {\textbf{Unlabelled}\\\textbf{batch $B_t$}\\[3pt] $n$ rows $x_{t,i}$\\[3pt] predictions\\$p,\ \hat{y},\ s$};

% --- input channel
\node[box=c3, text width=27mm, right=8mm of batch.north east, anchor=north west] (ks) {\textbf{KS tests}\\$M$ features + 3 PCA\\coords; Holm, $\alpha{=}.05$};
\node[box=c3, text width=24mm, right=5mm of ks] (alarm) {\textbf{Input alarm} $a_t$\\+ correlation\\statistic $\Delta_t$};
\node[box=c3, text width=24mm, right=5mm of alarm] (loc) {\textbf{Localization}\\rank by $D_{t,j}$;\\flag if $\tilde\pi_{t,j}\le\alpha$};

% --- impact channel
\node[box=c1, text width=27mm, right=8mm of batch.south east, anchor=south west] (q) {\textbf{Calibrated AC}\\$q_{t,i}=g(s_{t,i})$\\$\Ahat_t=\frac{1}{n}\sum_i q_{t,i}$};
\node[box=c1, text width=24mm, right=5mm of q] (pb) {\textbf{Poisson-binomial}\\$F_t$; 95\% plug-in\\interval};
\node[box=c1, text width=24mm, right=5mm of pb] (kap) {\textbf{Harm concern}\\$\concern_t=F_t(k_{\max})$\\flag if $\concern_t\ge .95$};

% --- lanes
\begin{scope}[on background layer]
\node[lane=c3, fit=(ks)(alarm)(loc)] (laneA) {};
\node[lane=c1, fit=(q)(pb)(kap)] (laneB) {};
\end{scope}
\node[lab=c3, anchor=south east] at (laneA.north east) {input channel (label-free)};
\node[lab=c1, anchor=north east] at (laneB.south east) {impact channel (label-free)};

% --- ledger, broker, verification, view
\node[box=c5, text width=27mm, right=9mm of loc.east, anchor=west, yshift=-2mm] (ledger) {\textbf{Evidence ledger}\\alarms, ranks, $\Ahat_t$,\\interval, $\concern_t$};
\node[box=c2, text width=27mm, below=4mm of ledger] (broker) {\textbf{Label broker}\\releases $y_b$ iff\\$b+d\le c$};
\node[box=c2, text width=27mm, below=4mm of broker] (verif) {\textbf{Verification}\\$A_b$, $H_b$, age $c-b$;\\else \emph{unresolved}};
\node[box=c4, text width=26mm, right=6mm of broker] (view) {\textbf{Aligned evidence}\\\textbf{view}\\shared batch axis,\\as-of clock $c$};

% --- arrows
\draw[arr] (fit.south) -- ++(0,-5mm) -| (batch.north);
\draw[arr] (batch.east |- ks.west) -- (ks.west);
\draw[arr] (ks) -- (alarm);
\draw[arr] (alarm) -- (loc);
\draw[arr] (batch.east |- q.west) -- (q.west);
\draw[arr] (q) -- (pb);
\draw[arr] (pb) -- (kap);
\draw[arr] (ref.south) -- ++(0,-4mm) -| (ks.north);
\draw[arr] (null.south) -- ++(0,-6mm) -| (alarm.north);
\draw[arr] (cal.south) -- ++(0,-2.5mm) -| ($(batch.east)+(4mm,0)$) |- ([yshift=2.5mm]q.west);
\draw[arr] (loc.east) -- (loc.east -| ledger.west);
\draw[arr] (kap.east) -| ([xshift=-4mm]ledger.south west) -- ([yshift=2mm]ledger.south west);
\draw[arr] (broker) -- (verif);
\draw[arr] (ledger.east) -| (view.north);
\draw[arr] (verif.east) -| (view.south);
\node[font=\scriptsize\itshape, text=c2!80!black, below=0.8mm of verif, text width=32mm, align=center] {ground truth never reaches the two channels};

% --- stage badges
\node[badge=cgray] at (fit.north west) {1};
\node[badge=c4] at (batch.north west) {2};
\node[badge=c3] at (ks.north west) {3};
\node[badge=c1] at (q.north west) {4};
\node[badge=c2] at (broker.north west) {5};
\node[badge=c4] at (view.north west) {6};
\end{tikzpicture}
}
\caption{\idm{} architecture. The classifier, scaler, PCA basis, isotonic
correctness map and reference sample are fitted on historical partitions and then
frozen (top). Each unlabelled batch feeds two independent channels. The input
channel applies per-feature and per-PCA-coordinate Kolmogorov--Smirnov tests with Holm
correction, an optional pairwise-correlation statistic and a feature ranking. The
impact channel computes calibrated average confidence, a plug-in Poisson-binomial
interval and a harm-concern score. Ground-truth labels are held by a release broker
and the labels of source batch $b$ reach the verification ledger only when
$b+d\le c$, where $c$ is the current clock. The aligned evidence view
reads the ledger as of a chosen clock $c$.}
\label{fig:arch}
\end{figure*}

\subsection{Input Channel: Detection}
\label{sec:detection}
Let $z=\phi(x)$ denote the standardized input. Let
$u=W^{\top}z\in\mathbb{R}^{3}$ be its projection on the first three principal
components, which are fitted on $\mathcal{D}_{\mathrm{tr}}$. For every original
feature and each of the three coordinates $u_1,u_2,u_3$, indexed jointly by
$j\in\{1,\dots,m\}$ with $m=M+3$, a two-sample
Kolmogorov--Smirnov (KS) statistic~\cite{massey1951ks} compares $B_t$ with a frozen
reference sample $R$ of $1{,}000$ historical rows:
\begin{equation}
D_{t,j}=\sup_{v}\bigl|\hat{F}_{B_t,j}(v)-\hat{F}_{R,j}(v)\bigr|.
\end{equation}
The $m$ resulting $p$-values $\pi_{t,j}$ are adjusted with Holm's step-down
procedure~\cite{holm1979}: with the $p$-values sorted in ascending order,
$\tilde{\pi}_{t,(i)}=\min\{1,\max_{l\le i}(m-l+1)\pi_{t,(l)}\}$.
The input alarm is $a_t=\mathbb{I}\{\min_j\tilde{\pi}_{t,j}\le\alpha\}$ with
$\alpha=0.05$. The PCA coordinates take the place of the autoencoder latent space
proposed in Phase~2 (Section~\ref{sec:changes}). They add a small amount of joint
evidence, but they do not reproduce a learned embedding detector such as
DriftLens~\cite{greco2025driftlens}.

Marginal tests cannot see a change that alters only feature dependence. An
optional \emph{dependence statistic} therefore tracks the largest absolute change in
pairwise Pearson correlation:
\begin{equation}
\Delta_t=\max_{i<j}\bigl|\hat{\rho}_{B_t}(i,j)-\hat{\rho}_{R}(i,j)\bigr|,
\label{eq:dep}
\end{equation}
where $\hat{\rho}_{B}(i,j)$ is the sample Pearson correlation of features $i$ and $j$
in sample $B$, and alarms when $\Delta_t>\tau_\Delta$. The threshold $\tau_\Delta$ is the empirical
$95$th percentile of $\Delta$ over stationary development batches. The arg-max pair
is reported as dependence evidence. It is descriptive, not causal.

\subsection{Input Channel: Localization}
Features are ranked by $D_{t,j}$ in descending order. A feature is \emph{flagged}
only when $\tilde{\pi}_{t,j}\le\alpha$. When the set $S$ of changed features is known,
which is the case in synthetic marginal shifts, localization is scored by
$\mathrm{P@1}=\mathbb{I}\{r_1\in S\}$ and
$\mathrm{R@3}=|\{r_1,r_2,r_3\}\cap S|/|S|$, where $r_k$ is the $k$-th ranked
feature. Both scores are undefined when $S=\emptyset$, as under conditional-only or
dependence-only change. The ranking is a description of distributional change and
does not identify a cause.

\subsection{Impact Channel: Calibrated Accuracy Estimation}
\label{sec:impact}
The impact channel follows the confidence-based performance estimators studied by
Kivim\"aki et al.~\cite{kivimaki2025confidence} and the average-confidence and ATC
literature~\cite{garg2022atc,guillory2021predicting}. It does not claim a new
estimator. Its contribution is to run the estimator behind a release broker and to
score it against verified harm. The raw confidence $s$ of a classifier is often
miscalibrated~\cite{guo2017calibration}, so \idm{} fits a monotone correctness map
$g:[0.5,1]\rightarrow[0,1]$ on $\mathcal{D}_{\mathrm{cal}}$ by isotonic
regression~\cite{zadrozny2002transforming}:
\begin{equation}
g=\arg\min_{\psi\ \text{non-decreasing}}\ \sum_{i\in\mathcal{D}_{\mathrm{cal}}}
\bigl(\mathbb{I}\{\hat{y}_i=y_i\}-\psi(s_i)\bigr)^{2}.
\label{eq:iso}
\end{equation}
Equation~\eqref{eq:iso} is the only learning objective in the monitor. The frozen
classifier is not retrained. Each prediction receives a calibrated correctness
probability $q_{t,i}=g(s_{t,i})$. The batch accuracy estimate is the calibrated
average confidence (AC)
\begin{equation}
\Ahat_t=\frac{1}{n}\sum_{i=1}^{n}q_{t,i}.
\label{eq:ac}
\end{equation}
An estimator is scored by $|\Ahat_t-A_t|$ once $A_t$ has been verified.

\subsection{Impact Channel: Interval and Harm Concern}
Treating the $q_{t,i}$ as known, independent Bernoulli success probabilities gives
a plug-in Poisson-binomial model for the number of correct predictions,
$K_t=\sum_iZ_{t,i}$ with $Z_{t,i}\sim\mathrm{Bernoulli}(q_{t,i})$. Its probability
mass function is computed exactly by the recursion
$P_i(k)=(1-q_{t,i})P_{i-1}(k)+q_{t,i}P_{i-1}(k-1)$~\cite{hong2013poibin}.
With $F_t$ the resulting distribution function, the $95\%$ plug-in interval for
$A_t$ is $[F_t^{-1}(0.025)/n,\;F_t^{-1}(0.975)/n]$. Harm~\eqref{eq:harm} is
evaluated in exact rational arithmetic. The largest harmful count is
$k_{\max}=\lceil n(\Aref-\delta)\rceil-1$, and the \emph{harm-concern score} is
\begin{equation}
\concern_t=\Pr\nolimits_{\mathrm{plug\text{-}in}}(K_t\le k_{\max})=F_t(k_{\max}),
\label{eq:concern}
\end{equation}
with a concern flag raised when $\concern_t\ge 0.95$. The interval is a model-based
prediction interval for realised batch accuracy. It is not a distribution-free
confidence interval. It ignores uncertainty in $g$ and in $\Aref$, and it gives no
guarantee under arbitrary shift. The input alarm is \emph{not} required before
$\Ahat_t$ and $\concern_t$ are computed or displayed.

\subsection{Label-Release Broker and Delayed Verification}
\label{sec:broker}
Online feature batches carry no label field. Labels are held by a broker that
releases the labels of batch $b$ at clock $c$ if and only if
\begin{equation}
b+d\le c.
\label{eq:release}
\end{equation}
Each clock has two phases. Predictions and channel outputs for $B_c$ are persisted
first, and released labels are then joined to the verification ledger. Any
component that \emph{fits} on released labels at forecast time $t$, such as a
recalibrator, a forecaster or an adaptive model, uses the stricter rule $b+d<t$.
Labels from the batch being forecast therefore never enter its own estimate. The
ledger records verified accuracy, harm and the age of the evidence. Cells that are
still unreleased are labelled ``unresolved''. They are never zero-filled or
back-dated.

\subsection{Aligned Evidence View}
The prototype view, an offline HTML replay, is controlled by task, seed,
condition, model, delay, as-of clock and batch. It aligns five panels on one batch
axis: the prospective estimate (calibrated AC with its interval, raw AC, an ATC-style
estimate~\cite{garg2022atc} and the concern badge); scheduled verification
(accuracy, macro-F1, verified harm and evidence age); a calibration diagnostic
available after release; the input evidence (alarms and the minimum Holm-adjusted
$p$); and a table of the same stream as of the selected clock. Its consistency with
the release rule was checked automatically (Section~\ref{sec:setup-impl}). No
analyst study was conducted.

\begin{algorithm}[t]
\caption{\idm{} monitoring loop for one stream}
\label{alg:loop}
\begin{algorithmic}[1]
\Require frozen $f,\phi,W,g,R,\Aref$; delay $d$; level $\alpha$; margin $\delta$
\For{$c=1,2,\dots,T+d$}
  \If{$c\le T$} \Comment{phase 0: unlabelled evidence}
    \State receive $B_c$; compute $p,\hat{y},s$; $q\gets g(s)$
    \State $D_{c,\cdot},\pi_{c,\cdot}\gets$ KS tests on $M$ features and 3 PCA coordinates
    \State $\tilde{\pi}\gets\textsc{Holm}(\pi_{c,\cdot})$; $a_c\gets\mathbb{I}\{\min\tilde{\pi}\le\alpha\}$
    \State rank features by $D_{c,j}$; flag $j$ if $\tilde{\pi}_{c,j}\le\alpha$
    \State $\Ahat_c\gets\mathrm{mean}(q)$; $F_c\gets\textsc{PoissonBinomial}(q)$
    \State $k_{\max}\gets\lceil n(\Aref-\delta)\rceil-1$; $\concern_c\gets F_c(k_{\max})$
    \State persist $a_c$, ranks, $\Ahat_c$, interval, $\concern_c$
  \EndIf
  \For{each $b$ with $b+d=c$} \Comment{phase 1: release}
    \State broker reveals $y_{b,\cdot}$; ledger stores $A_b$, $H_b$, age $c-b$
  \EndFor
\EndFor
\end{algorithmic}
\end{algorithm}

Algorithm~\ref{alg:loop} summarizes one monitoring stream. Its cost is dominated by
the $m$ KS tests, $O(mn\log n)$ per batch, and by the Poisson-binomial recursion,
which is $O(n^{2})$ per batch for $n=500$.

\subsection{Changes After Phase~2}
\label{sec:changes}
No instructor feedback on the Phase~2 proposal was received. The changes below come
from the authors' own audit of the proposal against what could be identified and
implemented. (i) The autoencoder latent anchor was replaced by a three-component PCA
fitted on training inputs. (ii) The proposed causal feature graph was not
implemented, and the feature evidence is reported as descriptive ranking only.
(iii) The forecast horizon was reduced to the current batch ($h=0$) in the core
pipeline. Prospective harm forecasting was moved to a separate study
(Section~\ref{sec:res-forecast}). (iv) The proposed analyst study was not attempted,
because no participant pool or ethics route existed within the course period; this
decision was recorded before any confirmatory outcome was seen. (v) Three additions
were made to close weaknesses found during implementation: the label-release broker,
exact rational harm arithmetic, and the dependence statistic~\eqref{eq:dep}.
